% AI-produced by GPT-6 Astra (OpenAI, ChatGPT Pro mode). No priority claim.
\section{The complete predimension family}
\label{sec:D}
Throughout this section $k\geq2$ and $a,b\geq1$ are integers,
$\delta(U)=a|U|-b|E(U)|$, and $\mathcal C(k,a,b)$ has the objects
and strong embeddings specified in the question. The zero ring is exact
for every category; every assertion of non-exactness below concerns a
nonzero unital ring. We allow noncommutative coefficient rings in the
explicit free-module obstructions, using left modules.

\begin{proposition}[Proved: closure calculus and functorial amalgamation]
\label{D:closure}
For every $(k,a,b)$ the class is a strong amalgamation class with
countably many isomorphism types. It satisfies (Rig), (SA), (FA) and
(FA$\varnothing$). Intersections of finite strong subsets are strong,
and restriction and composition preserve strong embeddings.
\end{proposition}
\begin{proof}
For subsets of an ambient finite hypergraph, edge counting gives
\[
 \delta(U)+\delta(V)\geq\delta(U\cup V)+\delta(U\cap V).
\]
If $A\leq_d B$ and $X\subseteq B$, then, for
$A\cap X\subsetneq U\subseteq X$,
\[
 \delta(U)-\delta(A\cap X)
 \geq\delta(U\cup A)-\delta(A)>0.
\]
Thus $A\cap X\leq_d X$. If $A\leq_d B\leq_d C$ and
$A\subsetneq U\subseteq C$, either $U\subseteq B$, or restriction
and submodularity give
$\delta(U)>\delta(U\cap B)\geq\delta(A)$. This proves
transitivity and then closure under intersections.

Let $A\leq_d B,C$, with the empty base permitted, and form the
hypergraph free amalgam $D=B\cup_A C$. For each $U\subseteq D$,
\[
 \delta_D(U)=\delta_B(U\cap B)+\delta_C(U\cap C)
             -\delta_A(U\cap A).
\]
If $U$ meets $C\setminus A$, the difference of the last two terms is
strictly positive, by restriction of closedness. Otherwise $U$ is a
subset of $B$. This proves positivity. The same equation, for
$B\subsetneq U$, proves $B\leq_d D$, and similarly for $C$.
The two images meet exactly in the base.

To make the operation functorial in $B$, in each gap of the ordered
base put the new $B$-points before the new $C$-points. A strong
embedding $B\to B'$ over $A$, together with the identity of $C$,
then induces an increasing induced embedding
$B\cup_A C\to B'\cup_A C$. Its source is strong: the target is
also the free amalgam of $B'$ and $B\cup_A C$ over $B$.
Composition is literal composition on tagged copies; transport to
any chosen skeleton gives the stated functor with its natural
inclusions. With empty base this is concatenation of the two orders.
Finite ordered structures are rigid, and a finite relational language
gives countably many finite isomorphism types. The free amalgams prove
joint embedding and amalgamation.
\end{proof}

\begin{theorem}[Proved: totality and open stabilizers for all parameters]
\label{D:total}
For every $(k,a,b)$, the limit $M$ satisfies (St) and (Op), and
$\mathcal C(k,a,b)$ is total. Consequently it is quasi-filtered.
\end{theorem}
\begin{proof}
We check the entire parameter-dependent step in T~A.30--A.34.
Let $B,C\leq_d D$, $F=B\cap C$, and let $d_0,d_1:D\to E$ be
strong embeddings agreeing on $F$. Choose an even integer $s=2N\geq8$.
Arrange one copy of $E$ and $2N-1$ copies $D_i$ of $D$ on a circle.
Identify consecutive copies over $B,C,B,C,\ldots$, using $d_0$ and
$d_1$ at the two ends. Put in exactly the hyperedges from the factors.
Every $F$-vertex has full-circle visibility. Every other vertex is
visible in at most three consecutive factors: away from $E$ it
crosses only one overlap, since it cannot belong to both
$B\setminus F$ and $C\setminus F$; at $E$ the two adjacent
occurrences can meet. Thus the factors inject and the overlaps are
exactly the specified ones.

An edge has $k$, rather than three, vertices, but this changes no
argument: if all its vertices occur in two factors, every one of their
visibility intervals contains the unique shortest path between those
factors. An edge therefore propagates through induced overlaps along
that path. The intersection of the vertex visibility intervals is an
interval, or the full circle. Consequently every factor is induced
and every edge identification is generated by consecutive overlaps.

Put $m=\max(|B|,|C|)$ and require
\[
                       s>3\bigl(\delta(E)+am\bigr).
\]
Cut one overlap to obtain a linear free amalgam $L$. Its factors are
strong. Restoring the cut identifies at most $m$ vertex pairs. If
$q:L\to Z$ is the quotient and $\widetilde U=q^{-1}U$, the exact
predimension loss is $a d_v-b d_e$, where $d_v\leq m$ and
$d_e\geq0$. Hence
\[
                      \delta_Z(U)\geq\delta_L(\widetilde U)-am.
\]
If the proper visibility intervals of $U$ miss an index, cut there:
all vertices and edges of $U$ lift uniquely, so positivity and strict
relative positivity over any contained factor follow from $L$.
Otherwise these intervals cover the circle. In the linear amalgam,
let $J$ consist of stages at which a non-$F$ vertex first appears.
Each such stage can account for at most three consecutive indices;
thus $3|J|\geq s$. At each stage in $J$, restriction of the incoming
strong base and submodularity show that the predimension increment
is a positive integer, hence at least one. Other increments are
nonnegative. Therefore
\[
 \delta_Z(U)\geq\lceil s/3\rceil-am>\delta(E)\geq\delta(D).
\]
This proves both positivity and closedness of every factor.

There is also an actual compatible order. Embed $E$ in $\mathbb R$,
write $p_t=(1-t)d_0+td_1$ on $D$, and let $\Delta>0$ bound below
all adjacent gaps of the $p_t$. For sufficiently large $N$, the maps
$p_{j/N}$ and the intermediate maps which retain the $B$-coordinates
at $p_{j/N}$ while moving all other coordinates to $p_{(j+1)/N}$
are strictly increasing: if
$L_0=\max_x|d_1(x)-d_0(x)|$, it suffices that
$2L_0/N<\Delta$. The intermediate maps agree successively over
$B$ and $C$, because the coordinates on $F$ do not change.
They descend to $Z\to\mathbb R$. Ordering by these values, with
arbitrary tie-breaking, preserves every factor order.

This finite bridge proves (St) algebraically. Indeed for
$g\in G_{B\cap C}$, take finite strong $D$ containing $B\cup C$
and finite strong $E$ containing $D\cup gD$, construct the bridge,
and embed it into $M$ over $E$. Extend each strong $D$-copy to an
automorphism $g_i$, with $g_0=1$, $g_{2N}=g$. Consecutive
$g_i^{-1}g_{i+1}$ fix alternately $B$ and $C$. Their product is $g$.
The reverse inclusion in (St) is immediate.

For an open $U\leq G$, choose a finite strong $B$ of least cardinality
with $G_B\leq U$. For $g\in U$, (St) gives
$G_{B\cap gB}\leq U$, whence $gB=B$ by minimality.
Its finite linear order makes $g|_B$ the identity. Thus $U=G_B$;
the empty support is permitted. This proves (Op).

For completeness, the remaining all-action argument has no
parameter-specific step. Free amalgamation shows
$\operatorname{Fix}_M(G_B)=B$: duplicate a finite strong extension
containing any proposed extra fixed point over $B$, and use homogeneity.
For every action $H$, condition I makes ``equal after an arrow'' an
equivalence relation and gives the injective-transition quotient
$H^{\rm inj}$. The maps
\[
 H\longrightarrow H^{\rm inj}\longrightarrow T_S,
 \qquad
 S=\mathop{\rm colim}_{B\subseteq M\ {
m strong}}H^{\rm inj}(B),
 \quad T_S(B)=S^{G_B},
\]
induce weak equivalences of categories of elements. To check this
rather than assume it, fix a target element in either comparison.
Its comma category is, after the slice comparison in the first case,
a nonempty full cosieve of $P/\mathcal C(k,a,b)$. The functor
$X\mapsto X\cup_P Z$, for any one object $Z$ of that cosieve,
has natural transformations from the identity and the constant $Z$,
so contracts it. Quillen's Theorem A applies twice. The action on
$S$ is $g\,i_B(h)=i_{gB}(H^{\rm inj}(g|_B)h)$ and is smooth.
By (Op) its orbits are $G/G_F$. For $F\ne\varnothing$,
$T_{G/G_F}(B)$ is naturally the set of strong embeddings $F\to B$,
using the fixed-point identity; its category of elements has an
initial object. For $F=\varnothing$ it is constant one; disjoint
ordered amalgamation contracts the category itself. Every action
component is therefore weakly contractible. This verifies the steps
of T~A.21 and A.30--A.34 used here, including the empty-base step.
\end{proof}

\begin{lemma}[Proved: a finite atom for every nontrivial sparse range]
\label{D:atom}
If $b<(k-1)a$, there is a finite ordered hypergraph $W$ on $n\geq3$
vertices such that
\[
 a<d:=\delta(W)\leq an/2,
 \qquad \delta(U)\geq d\quad(2\leq |U|<n).
 \tag{D1}\label{eq:D-atom}
\]
In particular $W$ is positive, every singleton is strong in $W$, and
its only proper nonempty strong subsets are singletons.
\end{lemma}
\begin{proof}
We give four explicit constructions, exhausting the range.

If $k=2$ and $a/2\leq b<a$, put $h=a-b$ and
$n=\lfloor a/h\rfloor+1\geq3$, and let $W$ be the cycle on
$n$ vertices. Then $d=nh\in(a,a+h]$ and $h\leq a/2$.
Every proper induced subgraph on $s\geq2$ vertices is a forest,
so its predimension is at least
$as-b(s-1)=hs+b\geq2a-b=a+h\geq d$.

Suppose $k\geq3$ and $b\leq a$. Set
$f(s)=as-b\binom{s}{k}$ and choose the least $n\geq4$ with
$f(n)\leq2a$. Such an $n$ exists. The consecutive differences of
$f$ are nonincreasing; also $f(2)=2a$, and $f(3)\geq2a$.
It follows that $f(s)\geq2a$ for $2\leq s<n$.
Choose any
\[
                      e=\left\lceil\frac{an-2a}{b}\right\rceil
\]
distinct $k$-subsets of the $n$ vertices as edges. Here $e\geq1$,
and $e\leq\binom nk$ because $f(n)\leq2a$. Now
$2a-b<d=an-be\leq2a$, so $d>a$ even when $b=a$.
Every proper $s$-vertex subset has predimension at least
$f(s)\geq2a\geq d$. Finally $n\geq4$ gives $d\leq an/2$.

For $k=2$ and $b<a/2$, repeat the preceding construction with
$L=2a-b$ in place of $2a$: take the least $n\geq4$ with
$f(n)\leq L$ and $e=\lceil(an-L)/b\rceil$ edges.
Here $f(2)=L$ and $f(3)>L$, so all proper subsets of size at least
two have predimension at least $L$. Moreover
$L-b=2a-2b>a$, and $a<d\leq L<2a\leq an/2$.

It remains to treat $b>a$. Put
\[
 q=\lceil b/a\rceil-1,\quad h=b-qa,\quad c=k-q,
 \qquad
 r=\left\lfloor\frac{(c-1)a-1}{h}\right\rfloor.
\]
Then $q\geq1$, $0<h\leq a$, $c\geq2$, and $r\geq1$.
Take a common core of $c$ vertices and $r$ disjoint petals of
$q$ vertices each; its $r$ hyperedges are core plus one petal.
Thus $n=c+rq$ and $d=ac-rh\in(a,a+h]$.
If a subset misses a core point, it has no edges, so its
predimension is at least $2a$ when its size is at least two.
A proper subset containing the core has $t<r$ full petals and
$s\geq0$ additional petal vertices; its predimension is
$ac-th+as>d$. If $n\geq4$, then $d\leq2a\leq an/2$.
The only possible $n=3$ has $c=2,q=r=1$; the maximality of $r$
then gives $2h\geq a$, whence $d=2a-h\leq3a/2$.

Any linear order may be placed on these structures. The final
closedness statements follow directly from (D1): every proper
extension of a singleton has larger predimension, whereas a proper
subset with at least two vertices has the full $W$ as an extension
of no larger predimension.
\end{proof}

\begin{theorem}[Proved: an all-ring obstruction throughout the sparse range]
\label{D:nonexact}
For $b<(k-1)a$ and every nonzero unital ring $R$,
$\mathcal C(k,a,b)$ is not exact over $R$. More precisely, for the
$W$ of Lemma~\ref{D:atom}, its singleton maps
$v_1,\ldots,v_n:V\to W$ have no affine equalizer over $R$.
Writing $P_T^R=R[\mathcal C(k,a,b)(T,-)]$, the natural map
\[
 D_W:P_W^R\longrightarrow(P_V^R)^{n-1},\qquad
 [g]\longmapsto([gv_i]-[gv_1])_{i=2}^n
 \tag{D2}\label{eq:D-column}
\]
is objectwise injective and
$T_1(\operatorname{coker}D_W)\cong R$.
\end{theorem}
\begin{proof}
Two distinct strong copies of $W$ in a target $X$ intersect in a
strong proper subset of either copy: apply restriction from
Proposition~\ref{D:closure} to the two strong images. Hence their
intersection has at most one vertex.
Distinct arrows have distinct images because $W$ is ordered.
Their sets of original $k$-edges are therefore disjoint.
For any $r>0$ distinct such copies, let $v$ be the number of vertices
in their union. Positivity of the induced union gives
\[
 av>b\,r|E(W)|=r(an-d)\geq\tfrac12anr.
\]
If every union vertex belonged to at least two copies, incidence
counting would give $nr\geq2v$, a contradiction. Thus some vertex
is private to exactly one copy.

Apply this observation to the nonzero support of a coefficient
vector in $P_W^R(X)$. At a private vertex, one component of its
image under (D2) has exactly the coefficient of that copy or its
negative. This is nonzero in every nonzero coefficient ring.
Thus $D_W(X)$ is injective. This is a natural map, since arrows act
by postcomposition on both free modules.

A covariant representable free module is projective, and its colimit
is $R$: its category of elements has initial object the identity.
On colimits each $v_i$ induces the identity of $R$, so
$\operatorname{colim}D_W=0:R\to R^{n-1}$.
The projective presentation
$0\to P_W^R\to(P_V^R)^{n-1}\to Q_R\to0$ therefore gives
$T_1(Q_R)=\ker(0)=R$. In particular an affine equalizer would
be an element of $\ker D_W(X)$ of augmentation one, impossible.
\end{proof}

\begin{theorem}[Proved: the complementary range is Ramsey and exact]
\label{D:high}
If $b\geq(k-1)a$, then $\mathcal C(k,a,b)$ is Ramsey, is exact
over every unital ring, and its ordered limit has extremely amenable
automorphism group. If $b\geq ka$, this category is precisely the
category of nonempty finite linear orders and increasing injections.
\end{theorem}
\begin{proof}
First assume $(k-1)a\leq b<ka$, and put
$\eta=b-(k-1)a\geq0$. A Berge cycle with $t$ edges has a union
of at most $(k-1)t$ vertices; its predimension is at most
$-\eta t\leq0$. Positivity forbids it, including a cycle of length
two arising when two edges meet in two vertices. Thus the incidence
graph is a forest. A connected component with $r$ edges has
$1+(k-1)r$ vertices and predimension $a-r\eta>0$.
Conversely a disjoint union of such finite hypertrees is positive:
every induced subhypergraph is a forest, and each of its components
has no more edges than an original component.

A strong subset must contain every edge it meets. Indeed, if it
contains exactly $j\geq1$ points of an edge but not all of it, adjoining
the missing $k-j$ points changes predimension by at most
$a(k-j)-b\leq a(1-j)\leq0$, contrary to strict closedness.
By connectedness it must contain each component it meets. Conversely
a union of components is strong, because any proper extension adds
a nonempty positive subset in the other components.

Encode these structures by their hyperedge relation and a unary
set-valued function $F(x)$ equal to the whole connected component
of $x$. Substructures in this language are exactly unions of
components, and embeddings are exactly our strong embeddings. This
is an isomorphism of the finite categories, not a change of the
embedding relation. The functional class is closed under free
amalgamation: over unions of whole components, no component from
one side meets a new component from the other side. In particular
the bound $a-r\eta>0$ is preserved. The class of all free linear
orders on it is Ramsey by Evans--Hubi\v cka--Ne\v set\v ril,
\emph{Ramsey properties and extending partial automorphisms},
Theorem~1.3 \cite{D-EHN2021}, whose language explicitly permits
unary set-valued functions of unbounded finite image size.

We spell out the exactness consequence to specify its strength.
The finite-hom Ramsey-to-affine-equalizer argument of T~A.6, rechecked
in its source, applies to every finite parallel family: if the
augmentation ideal of its integral kernel were $d\mathbb Z$ with
$d\ne1$, a circle-valued functional on the image of its difference
matrix, extended by divisibility, would give a finite coloring for
which no required copy is monochromatic. Partitioning each of the
$\ell$ circle coordinates into $3\ell+1$ intervals makes the
contradiction explicit. Thus there is an integral affine equalizer
for every family. Isbell--Mitchell's criterion gives exactness over
$\mathbb Z$; transporting the same integral coefficients through
$\mathbb Z\to R$ gives it over every unital $R$. The strong-class
KPT theorem, EHN Theorem~2.13 \cite{D-EHN2019}, applies to this
rigid strong amalgamation class and proves extreme amenability.

Finally a single hyperedge has predimension $ka-b\leq0$ when
$b\geq ka$, so no hyperedge occurs. Every subset of an edgeless
object is strong, giving exactly the asserted order category; its
Ramsey property is the finite Ramsey theorem. The same conclusions
follow.
\end{proof}

\begin{theorem}[Proved: amenability classification]
\label{D:amenability}
The group $\operatorname{Aut}(M)$ for the ordered class
$\mathcal C(k,a,b)$ is amenable if and only if it is extremely
amenable, if and only if $b\geq(k-1)a$. The same construction below
proves non-amenability for the unordered limit when $b<(k-1)a$.
\end{theorem}
\begin{proof}
Only the negative range remains. Put $h=\lceil a/b\rceil$.
Let $\Omega$ be the compact space of assignments of each hyperedge
$e$ of $M$ to one of its vertices, with no vertex assigned more
than $h$ edges. It is nonempty: for any finite edge family $J$
with union $U$, positivity gives
$|J|\leq|E(U)|<(a/b)|U|\leq h|U|$, which is Hall's condition
for matching edges to $h$ copies of the vertices. The finite
assignments pass to an infinite assignment by compactness.
The conditions defining $\Omega$ are closed in
$\prod_{e\in E(M)}e$. Permuting finite coordinates proves that the
natural $G$-action is continuous.

The single-edge structure $W_0$ has all singleton subsets strong
because $ka-b>a$. Fix a strong copy $e$ in $M$. If a
$G$-invariant probability measure existed on $\Omega$, some event
``$e$ is assigned to $x$'', $x\in e$, would have probability
$p>0$. The $G_x$-orbit of $e$ is infinite. Indeed freely amalgamate
arbitrarily many copies of $W_0$ over $x$, amalgamate the orders,
embed the result strongly into $M$, and extend the copy
isomorphisms by homogeneity. For distinct orbit edges
$e_1,\ldots,e_N$, invariance makes each assignment-to-$x$ event
have probability $p$, while the sum of their indicators is at most
$h$. Therefore $Np\leq h$ for every $N$, a contradiction.
This is the hypergraph version of the opened proof of EHN
Theorem~3.8; it does not invoke an invalid assertion about the shadow
graph. A nonempty compact $G$-space without invariant probability
measure proves non-amenability. Positive-range extreme amenability
was proved in Theorem~\ref{D:high}.
\end{proof}

\begin{proposition}[Proved and Refuted: the shadow-graph issue]
\label{D:shadow}
For $\mathcal C_\delta$, the graph joining vertices of a common
hyperedge is $2$-sparse, and the two stabilizer orbits required by
EHN Theorem~3.8 are infinite. More generally its $k$-uniform analogue
is sparse when $b>a(k-2)$. If $b\leq a(k-2)$, there is no uniform
sparsity bound for the shadow graphs. In particular EHN Remark~3.2,
as written, is false for arbitrary sparse relations.
\end{proposition}
\begin{proof}
For a finite vertex set $B$ in the ternary case, let $e=|E(B)|$.
Each shadow pair not accounted for by an internal edge has a
witnessing triple with exactly two points in $B$. Choose one per
pair, giving $t$ distinct triples and at most $t$ new vertices.
Positivity on their union implies
$3(e+t)<2(|B|+t)$, hence $t<2|B|-3e$. The number of shadow
pairs is at most $3e+t<2|B|$. The orbit assertion follows from
free amalgamation of strong single edges over either endpoint,
as in Theorem~\ref{D:amenability}.

In general choose $t$ distinct external witnesses covering all
remaining pairs. They add at most $(k-2)t$ points. Thus
\[
 [b-a(k-2)]t<a|B|-b|E(B)|.
\]
When the left coefficient is positive, both $t$ and $|E(B)|$ are
bounded linearly in $|B|$, and every witness accounts for at most
$\binom{k}{2}$ pairs. This proves sparsity.

For the converse, take an arbitrary set $B$ and, for every pair
of its vertices, add $k-2$ private new vertices and one hyperedge
consisting of that pair and its private vertices. Every subset $U$
has
\[
 \delta(U)\geq a|U\cap B|+
          [a(k-2)-b]r>0
\]
when $U$ contains an edge, where $r$ is the number of complete
edges in $U$; subsets with no edge are positive as well.
Its shadow on $B$ is the complete graph. Taking arbitrarily large
$B$ disproves a uniform bound.

For a concrete counterexample to the published remark itself, take
eight base vertices and one private vertex $z_{ij}$ for each of
its 28 pairs, with the 28 ordered relational tuples $(i,j,z_{ij})$,
$i<j$. This ternary relation is $1$-sparse on every subset, because
the tuple-to-private-vertex map is injective. Its shadow induces
$K_8$, with $28>\binom{1\cdot3}{2}\cdot8=24$ edges.
The finite checks are in \texttt{code/D\_certificates.py}.
We also verify the particular potential construction in Lead~4(c).
For a finite probability vector $(t_g)$ on strong triples, put
$\mu_i(v)=\sum_{gv_i=v}t_g$ and
$w(v)=\max_{g(W)\ni v}t_g$. Applied to the threshold supports
$\{g:t_g>s\}$, the inequality $2|U|>3|S|$ above gives, by
integrating finite step functions,
$\sum_v w(v)\geq3/2$. Since
\[
 w(v)\leq\max_i\mu_i(v)\leq
 \mu_1(v)+(\mu_2(v)-\mu_1(v))_+
          +(\mu_3(v)-\mu_1(v))_+,
\]
and each marginal has mass one, we recover exactly T~A.18:
$\|\mu_2-\mu_1\|_1+\|\mu_3-\mu_1\|_1\geq1$.
Finite-dimensional separation gives functions
$\phi_1,\phi_2$ on the strong points of any finite strong target,
with values in $[-1,1]$, such that
\[
 \phi_1(gv_2)-\phi_1(gv_1)+
 \phi_2(gv_3)-\phi_2(gv_1)\geq1
\]
for every strong triple in that target. Indeed separate their
column convex hull from the open unit $\ell^1$ ball and normalize
the separating functional in the dual supremum norm.
Let $O$ be the single $G$-orbit of strong singleton points of $M$.
The set of all pairs in $([-1,1]^O)^2$ satisfying every such
inequality is nonempty by compactness: finitely many inequalities
lie in one finite strong target. It is a compact convex set with
a continuous coordinate-permuting $G$-action. A fixed pair would
have both functions constant on $O$, contradicting $0\geq1$.
This proves the stated direct potential route; the route through
(Re) uses the displayed quantitative inequality and Part~A.
\end{proof}

\begin{theorem}[Proved: no precompact Ramsey expansion in the negative range]
\label{D:no-precompact}
If $b<(k-1)a$, the ordered limit $M$ has no co-precompact extremely
amenable subgroup. Consequently it has no precompact reasonable
strong Ramsey expansion whose Fra\"isse limit is an expansion of
$M$ in the sense defined in EHN immediately after Remarks~2.19. Orientation expansions
with closures must not be described as precompact over this original
limit merely because each finite object has finitely many orientations.
\end{theorem}
\begin{proof}
Use the assignment space $\Omega$ from
Theorem~\ref{D:amenability}. An extremely amenable subgroup $H$
fixes an assignment. Orient each assigned hyperedge from its assigned
vertex toward its other $k-1$ vertices. The resulting directed graph
has out-degree at most $d_0=h(k-1)$, so every finite-radius forward
ball is finite.

Put $c=(k-1)a-b>0$ and choose $L\geq2$ with $Lc>a$.
A finite loose hypergraph tree is positive: a component with $r$
edges has predimension $a+rc$. Two vertices whose unique connecting
hyperedge path has length $t\geq L$ form a strong edgeless
2-point subset. To verify this, an induced subset containing both
vertices either keeps them in different components, in which case
every proper extension of the pair has predimension greater than
$2a$, or connects them, in which case it contains all $t$ path
edges and has predimension at least $a+tc>2a$.

For every height $m$, take a rooted loose tree in which each
non-leaf vertex has $h+1$ child hyperedges, each with $k-1$ new
children. Order it by depth, breaking ties arbitrarily, and embed
it strongly into $M$. In any assignment, some child hyperedge at
each non-leaf vertex is assigned to a child, since only $h$ edges
can be assigned to that vertex. Following such children gives a
path which, read upwards from its bottom, is a directed path
$x_0\to x_1\to\cdots\to x_m$, strictly decreasing in the chosen
linear order. For every $i\geq L$, the pair $(x_i,x_0)$ is a
strong copy of the same ordered edgeless 2-point structure.

If $H$ were co-precompact, it would have only finitely many orbits
on that one $G$-orbit of ordered pairs; say $t_0$. Thus for any
fixed endpoint $x_0$, the vertices $x_i$, $i\geq L$, lie in
at most $t_0$ orbits under $H_{x_0}$. This is incompatible with
arbitrarily long such directed paths. Indeed an orbit which meets
the forward ball of radius $j$ is contained in that ball, because
$H_{x_0}$ preserves the assignment. Starting with any initial
segment, once its finite ball is exceeded by the number of distinct
path vertices, a new orbit must occur. Iterating this observation
$t_0+1$ times yields a finite bound depending only on $d_0,L,t_0$
for the path length. Equivalently one may use the explicit bound
$B(j)=1+d_0+\cdots+d_0^j$ and indices
$j_1=L$, $j_{r+1}=B(j_r)+L+1$: the segment through $j_{r+1}$
contains a vertex outside the union of all orbits previously
represented through $j_r$. Choosing $m\geq j_{t_0+1}$ contradicts
the orbit bound.

This reproduces the specific finite-orbit argument of the opened
EHN Theorem~3.16(1), with loose hypergraph trees and the fixed order
included. The final Ramsey-expansion implication is the strong-class
KPT correspondence and the definition of co-precompactness.
\end{proof}

\begin{corollary}[Proved: the complete parameter table]
\label{D:table}
For every prime $p$, the following table gives all parameters.
Every category in the table satisfies (Rig), (SA), (FA),
(FA$\varnothing$), (St), and (Op).
\[
\begin{array}{c|ccccccc}
 \text{range}&\text{filtered}&\text{exact }\mathbb Z&
 \text{exact }\mathbb Q&\text{exact }\mathbb F_p&
 \text{total}&\text{Ramsey}&G\text{ amenable}\\ \hline
 b<(k-1)a&\text{no}&\text{no}&\text{no}&\text{no}&
 \text{yes}&\text{no}&\text{no}\\
 b\geq(k-1)a&\text{no}&\text{yes}&\text{yes}&\text{yes}&
 \text{yes}&\text{yes}&\text{yes}
\end{array}
\]
In both rows quasi-filteredness holds. Extreme amenability agrees
with amenability. No member of this family refutes A.3(i) or A.3(iv).
In particular $\mathcal C(2,2,1)$ and $\mathcal C(3,3,1)$ are total,
not exact over any nonzero ring, and have non-amenable automorphism
groups.
\end{corollary}
\begin{proof}
The preceding theorems give all entries except filteredness.
The two different embeddings of a singleton into an edgeless ordered
pair cannot become equal after any embedding, because every arrow
is injective. Thus none of these categories is filtered. In the
negative range non-Ramseyness also follows directly from the
non-exactness theorem and the Ramsey-to-exactness implication.
For the two highlighted examples, the atoms can be taken as the
triangle, respectively a 5-vertex ternary hypergraph with any nine
of its ten triples; the certificate gives them explicitly.
\end{proof}

\begin{proposition}[Proved: the five-edge certificate and its limits]
\label{D:five-edge}
Lead~8(e)'s displayed integral affine equalizer is correct whenever
$a>2b$. Its five-edge target is positive exactly when $a>b$, and
all its displayed hyperedges are strong exactly when $a>2b$.
This equalization of the initial triple does not imply exactness.
\end{proposition}
\begin{proof}
On ordered vertices $1<\cdots<5$ take edges
$123,234,345,124,245$. For
$t=[123]+[234]+[345]-[124]-[245]$, direct collection in each
position gives $tv_1=tv_2=tv_3=[3]$ and augmentation one.
The full object has predimension $5(a-b)$, so positivity requires
$a>b$. Conversely every subset has at most as many edges as
vertices, so that inequality suffices. A listed edge has
predimension $3a-b$. A four-vertex extension contains at most
three edges, and a five-vertex extension contains five, so its
relative predimension is at least $a-2b$, respectively $2a-4b$.
The four-vertex sets $1234$ and $2345$ attain the first bound,
and the whole object attains the second. Thus all the listed
edges are strong exactly when $a>2b$. The finite marginals and
boundary cases are checked by the program. Theorem~\ref{D:nonexact}
provides another family without an affine equalizer even in this
range. Also, if every vertex of the support belongs to exactly two
triples, its three equal position-marginals must all vanish at that
vertex, since at least one position is absent there. Thus all
marginals vanish, and their augmentation, which is the augmentation
of the coefficient vector, is zero. Such a support cannot give an
affine equalizer over any ring. Finally, decreasing $b/a$ preserves
positivity and the closedness of every fixed finite embedding, so
the parameter set where this particular triple can be equalized
is downward closed; this explains the threshold remark without
conflating it with the exactness classification.
\end{proof}

\begin{theorem}[Proved: the unordered class of $\mathcal C_\delta$ is not exact]
\label{D:unordered}
In the unordered class, the three singleton embeddings into one
hyperedge have the rational affine equalizer
$\frac16\sum_{\sigma\in S_3}[\sigma]$. Nevertheless its category
is not exact over any nonzero ring, including $\mathbb Q$.
\end{theorem}
\begin{proof}
The averaging formula is literal: each of the three original
vertices is carried to each target vertex by two permutations, so
all three marginals are one third of the vertex sum.

For the obstruction we give an explicit rigid replacement atom.
Let $J$ be the 12-vertex cubic graph with edge set
\[
\begin{split}
&01,07,0\,11,12,1\,11,23,2\,10,34,35,45,49,56,\\
&67,68,78,89,9\,10,10\,11.
\end{split}
\]
Here $uv$ denotes the unordered pair $\{u,v\}$, not a number.
The finite certificate enumerates every nontrivial edge cut and
every adjacency-preserving vertex permutation: the minimum cut
has size three and the only automorphism is the identity.
The enumeration algorithm extends a connected ordering of the
vertices, testing adjacency and non-adjacency against every
previously assigned vertex; it exhausts all permutations.

Take the 18 edges of $J$ as the vertices of a 3-uniform hypergraph
$W$. For each vertex of $J$ other than $0$, make its three incident
edge-tokens a hyperedge of $W$. Thus $W$ has 18 vertices and 11
hyperedges, and $\delta(W)=36-33=3$.
For a vertex subset $U$ of $W$ containing at least one hyperedge,
let $S$ be the set of all hyperedges contained in $U$, identified
with a subset of $V(J)\setminus\{0\}$. Their union $U_0$ has
$3|S|-|E_J(S)|$ vertices. It contains exactly the edges $S$: any
additional complete edge would be contained in $U$ already. Hence
its predimension is
\[
              3|S|-2|E_J(S)|=|\partial_J S|\geq3.
\]
The remaining vertices $U\setminus U_0$ contribute two each
and no further edges, so $\delta(U)\geq3$. A subset of size at least
two with no hyperedge has predimension at least four. Singletons
have predimension two. Thus only singletons are proper nonempty
strong subsets of $W$.

The hypergraph reconstructs $J$: its hyperedges are the 11
nondeleted graph vertices; degree-two edge-tokens join their two
incident hyperedges; its three degree-one tokens join to one new
vertex, recovering $0$. Hence a hypergraph automorphism extends
to an automorphism of $J$ fixing $0$, and is the identity. $W$
is rigid without imposing an order.

Distinct embeddings of this $W$ now have distinct images, so the
private-vertex proof of Theorem~\ref{D:nonexact} applies unchanged.
Indeed distinct strong images meet in at most one vertex, their
11 original triples are disjoint, and positivity forces a private
vertex because $\delta(W)=3\leq2\cdot18/2$. Its 18 singleton
maps therefore have no affine equalizer over any nonzero ring.
The same projective cokernel has $T_1\cong R$.
The program also checks all 262143 nonempty vertex subsets of $W$.
\end{proof}

\begin{proposition}[Proved and Refuted: the expansion spectral sequence]
\label{D:expansion}
Suppose $u:A^*\to\mathcal C_\delta$ is any Ramsey expansion
category with finite hom-sets and common successors, and let $R$
be a nonzero unital ring. Then $\operatorname{colim}_{A^*}$ is
exact and $u$ is not homotopy final. The proposed description of
$T_1(Q_R)$ as the image of first comma homology under the stated
spectral sequence is false: every term $E^2_{1,0}$ of that
spectral sequence is zero, so it cannot be $T_1(Q_R)\cong R$.
\end{proposition}
\begin{proof}
Ramsey gives integral affine equalizers, hence exactness over $R$.
If $u$ were homotopy final, for every $R$-module diagram $V$ on
$\mathcal C_\delta$ the natural comparison of derived colimits
\[
 \mathop{\rm Lcolim}_{A^*}u^*V\longrightarrow
 \mathop{\rm Lcolim}_{\mathcal C_\delta}V
\]
would be a quasi-isomorphism: the two-sided bar comparison has
contractible comma categories $X/u$ as its fibers. Exactness on
$A^*$ would force $T_1(V)=0$, contradicting (D2).

Here is the correctly typed spectral sequence. Its coefficient
is a diagram $F:A^*\to R\text{-}\mathrm{Mod}$, and
\[
 (L_q u_!F)(X)=H_q(u/X,F\circ\pi_X),
 \qquad
 E^2_{p,q}=T_p^{\mathcal C_\delta}(L_q u_!F)
 \Longrightarrow T_{p+q}^{A^*}(F).
 \tag{D3}\label{eq:D-Kan-ss}
\]
The comma objects are arrows $u(a^*)\to X$, and
$\pi_X:u/X\to A^*$ sends them to $a^*$. This is $u/X$, whereas
homotopy finality above uses $X/u$. The left Kan functor $u_!$
preserves projectives, because its right adjoint $u^*$ is exact;
therefore the projective-complex construction gives (D3).
It is a first-quadrant homological spectral sequence, with
$d_r:E^r_{p,q}\to E^r_{p-r,q+r-1}$. No differential enters or
leaves $(1,0)$ for $r\geq2$. Its abutment in total degree one is
zero, hence
\[
                       T_1^{\mathcal C_\delta}(u_!F)=0
                       \quad\text{for every }F.
\]
Consequently $Q_R$ is not isomorphic to $u_!F$ for any $F$.
Taking $F=u^*Q_R$ does not change this: the counit
$u_!u^*Q_R\to Q_R$ cannot be an isomorphism.
The degree-one five-term sequence instead relates
$T_2(u_!F)$ and $\operatorname{colim} L_1u_!F$ before the zero
abutment; it has no edge map producing the claimed nonzero
$T_1(Q_R)$. This corrects the third clause of FE-4 without
presupposing that a particular orientation construction is
precompact.
\end{proof}
